\chapter{Orthodromic Rivers: Geometry}
\label{ch:17}

The first terrestrial program to be treated is the family of freshwater corridors that the working vocabulary of this book calls orthodromic rivers. The term refers to engineered channels or conduits that follow, as nearly as terrain and the boundaries of the oceans permit, the arcs of great circles, which are the shortest paths between points on a sphere and are called orthodromes in the vocabulary of navigation. The concept belongs to a broader scheme set out in an earlier work by the author, the Orthodromic Infrastructure Blueprint, in which nine great circles, the mirror planes of the octahedral symmetry group, are used to organize a planetary network of infrastructure in a manner that is symmetric and free of privileged regions. The present book takes the hydraulic layer of that scheme as its subject, and it treats the geometry and the water separately, in this chapter and the next. It adopts as a working definition that an orthodromic river is a corridor for the conveyance of fresh water laid along an arc of one of these circles, subject to the modifications that the terrain and the law require, and it flags the definition as provisional pending the author's more detailed statement of the water system's parameters.

The attraction of great-circle routes is evident from the elementary geometry of the sphere, but the attraction needs to be examined and not assumed. On a sphere the shortest path between two points is the arc of the great circle through them, and the saving over other natural routes can be substantial: two points at latitude $60^\circ$ and separated by $90^\circ$ of longitude are $4{,}605$ km apart along the great circle, whereas the path along the parallel of latitude is $5{,}004$ km, longer by eight per cent. Shorter conduits cost less to build, lose less water to seepage and evaporation, and require less energy for pumping against friction. These advantages are real. They are, however, advantages of shortness, and the question whether shortness is the property to be minimized depends on the cost structure of the terrain. A route of minimum length may cross a mountain range which a longer route avoids, and the saving in length is more than lost in the cost of tunnelling or pumping. The correct generalization is the principle that governs optimal paths in any medium with variable cost, which is that the optimal route is a geodesic of the metric weighted by the cost of passage. A great circle is optimal only where that cost is uniform, and in all other cases it serves as the reference against which the penalty of the actual route is measured. This is the reading of the great-circle network that the book adopts: it is a reference configuration that supplies a lower bound on length and a symmetric organization of the globe, and the corridors that would actually be built would deviate from it to the extent that terrain, jurisdictions, and ecology demand.

The octahedral scheme has a combinatorial structure that can be set out exactly. The full symmetry group of the octahedron contains forty-eight elements and nine mirror planes, three of which are the coordinate planes and six of which are the planes bisecting pairs of coordinate axes. Each plane meets the sphere in a great circle, and the nine circles divide the sphere into forty-eight congruent spherical triangles with interior angles of $90^\circ$, $60^\circ$, and $45^\circ$. The vertices at which the circles meet are of three kinds: six at the ends of the coordinate axes, where four circles meet, eight at the directions of the cube's vertices, where three circles meet, and twelve at the midpoints of the edges of the octahedron, where two circles meet. Each circle passes through exactly eight of the twenty-six vertices and is therefore divided into eight arcs, so that there are seventy-two arcs in all. The arcs are of three lengths, corresponding to angular lengths of $35.26^\circ$, $45^\circ$, and $54.74^\circ$, twenty-four of each, which on a sphere of the radius of the Earth are about $3{,}921$ km, $5{,}004$ km, and $6{,}086$ km. The total length of the nine circles is $9\times2\pi R\approx3.6\times10^5$ km. These numbers fix the scale of the network that the book considers: twenty-six nodes joined by seventy-two arcs of a few thousand kilometers each, and a total length that is about nine times the circumference of the Earth.

The Blueprint itself distinguishes this concurrent, cubic-symmetric realization from the general case. In general position, nine great circles intersect pairwise in two antipodal points and so produce $2\binom{9}{2}=72$ vertices, $144$ edges, and $74$ faces, a count that satisfies $72-144+74=2$ and is the upper bound for any arrangement of nine circles; in the Blueprint's general formulation, five of the circles are fixed as reference orthodromes and four are inclined and flexible, so that the network can be adjusted to the geography without losing the protocol. The cubic-symmetric case, in which the nine circles are the mirror planes of full octahedral symmetry, is the special arrangement in which many circles meet at common points and the vertex count falls to $26$. The count has a familiar mnemonic, which the author has used in physical models of the arrangement: the $26$ nodes correspond to the $8$ corners, $12$ edges, and $6$ face centers of a cube, which is also the number of visible pieces of the puzzle cube. The author has built such models, including a globe with the circles drawn on it and a lattice of the Hoberman type hung over a globe, and these serve as physical checks on the geometry. The present book adopts the cubic-symmetric realization as its reference configuration, because it is exactly specified and its numbers can be computed, and it treats the flexible circles of the general case as the means by which the reference configuration would be adapted to terrain.

A very large share of such a network could not be built as a surface channel, because most of the surface is ocean. The land area of the Earth is about twenty-nine per cent of the whole, and on average the same proportion of a great circle lies over land, so that of the $3.6\times10^5$ km of the nine circles perhaps $1.0\times10^5$ km would pass over land, and the rest across water, where a corridor would have to be a submerged pipeline or would not exist at all. Moreover, the orientation of the reference configuration relative to the Earth's axis is a free parameter, as are the inclinations of the flexible circles in the general case, and different choices yield different land fractions and different intersections with mountain ranges, deserts, and the territories of the states that would have to participate. The selection of an orientation is a constrained optimization problem over the group of rotations, and its objective is not geometric alone. It must weigh the water demand of the regions traversed against the cost of traversal, the number of international borders crossed, and the sensitivity of the ecosystems along the route. The book states the problem and leaves its solution to the stage at which regional data are available. It treats the preliminary conclusion as one that follows from the geometry: the network, in its pure form, is not a plan for the whole planet but a plan for those arcs, or portions of arcs, that fall on land where freshwater demand and supply are mismatched, and its value must be judged arc by arc.

A second geometrical point is that freshwater conveyance is constrained by elevation in a way that has no counterpart in the flat geometry of navigation. A great circle takes no account of the relief it crosses. Water flows downhill without cost and must be lifted at cost, so the desirable corridor is one whose profile of elevation lets gravity do as much of the work as possible and requires lift only where the demand lies above the source. Because the energy of lift scales with the height gained, as the next chapter shows, a corridor that crosses a ridge of one kilometer must expend energy of the order of $2.7$ kWh per cubic meter to carry the water over it, and a route that follows contours and keeps a constant gentle gradient may spend none. The great-circle reference therefore must be supplemented by a second reference, the profile of elevation along the route, and an assessment of a corridor consists of the joint consideration of both.

The geometry also bears on questions of governance that are treated in the next chapter. A corridor of thousands of kilometers crosses many jurisdictions, and its symmetric layout is not a source of legitimacy. A scheme that treats the planet as a uniform surface to be organized by a single geometric rule has the character of the abstract plans against which the literature on large-scale state projects has warned, in which a legible and elegant design imposed from above fails to take account of the local knowledge and the existing arrangements of the people affected. The response is not to abandon geometry, which supplies lower bounds and avoids arbitrary favoritism, but to treat it as a language in which the proposals are stated and compared and not as a source of authority. Where a route deviates from the great circle in order to avoid a sacred site, a settlement, or a wetland, the deviation is a reasonable decision, and its length penalty is simply a price to be reported.

\section*{Formalization}

The spherical geometry is fixed first. Let the Earth be modelled as a sphere of radius $R=6371$ km; the departure of the true figure from a sphere is a flattening of about $1/298$, which alters distances by at most a few tenths of one per cent and is neglected here. The angular distance $\theta$ between points with latitudes $\varphi_1,\varphi_2$ and longitude difference $\Delta\lambda$ satisfies $\cos\theta=\sin\varphi_1\sin\varphi_2+\cos\varphi_1\cos\varphi_2\cos\Delta\lambda$, and the orthodromic distance is $R\theta$.

\begin{proposition}
Let $\Gamma$ be the group generated by reflections in the nine planes with normals $e_1,e_2,e_3$ and $(e_i\pm e_j)/\sqrt2$, $i<j$. Then $|\Gamma|=48$; the nine great circles divide the sphere into $48$ congruent spherical triangles with angles $\pi/2,\pi/3,\pi/4$, each of area $4\pi R^2/48\approx1.06\times10^7\ \mathrm{km^2}$; the vertices number $26$ (of valencies $4$, $3$, $2$, in numbers $6$, $8$, $12$); each circle contains $8$ vertices; and the arcs number $72$, of which $24$ have angular length $\arccos(1/\sqrt2)=45^\circ$, $24$ have $\arccos(\sqrt{2/3})\approx35.26^\circ$, and $24$ have $\arccos(1/\sqrt3)\approx54.74^\circ$.
\end{proposition}

\begin{proof}
The group generated by these reflections is the full octahedral group $O_h$, of order $48$. A spherical triangle with angles $\pi/2,\pi/3,\pi/4$ has area $R^2(\pi/2+\pi/3+\pi/4-\pi)=R^2\pi/12=4\pi R^2/48$ by Girard's theorem, so $48$ such triangles tile the sphere. The vertices are the points of the sphere fixed by a nontrivial rotation subgroup: the $6$ axis directions $(\pm1,0,0)$ and permutations, the $8$ directions $(\pm1,\pm1,\pm1)/\sqrt3$, and the $12$ directions $(\pm1,\pm1,0)/\sqrt2$ and permutations, giving $26$. Counting incidences, a circle with normal $n$ contains the vertices orthogonal to $n$, and direct enumeration gives $8$ for each of the nine normals. Each circle is thus divided into $8$ arcs, giving $72$. The Euler relation $V-E+F=26-72+48=2$ is satisfied. The three angular lengths are the angles between adjacent vertices of the three types: between an axis and an edge-midpoint, $\arccos(1/\sqrt2)=45^\circ$; between a cube-vertex and an edge-midpoint, $\arccos(\sqrt{2/3})$; between an axis and a cube-vertex, $\arccos(1/\sqrt3)$. Each type of edge occurs $24$ times, and the angles sum to $24(45+35.26+54.74)=3240^\circ=9\times360^\circ$, as required for nine full circles.
\end{proof}

The sums have been verified by direct enumeration of the vertices and edges. On the sphere above, the arc lengths are $5003.8$, $3921.2$, and $6086.3$ km, and the total length of the nine circles is $9\cdot2\pi R=360{,}272$ km.

The second result states the sense in which a great circle is optimal. Let a corridor be a curve $\gamma$ on the sphere joining two points, and let the cost of construction and operation per unit length at a point $x$ be $c(x)>0$, a continuous function. Define the cost $C(\gamma)=\int_\gamma c(x)\,ds$.

\begin{proposition}
A minimizer of $C$ among curves with fixed endpoints is a geodesic of the conformally rescaled metric $c(x)^2\,g$, where $g$ is the round metric. If $c$ is constant, the minimizers are arcs of great circles, and $C=c\,R\theta$. In the plane approximation, at an interface across which $c$ jumps from $c_1$ to $c_2$, the minimizer satisfies the refraction relation $c_1\sin\alpha_1=c_2\sin\alpha_2$, where $\alpha_i$ is the angle between the path and the normal to the interface on side $i$.
\end{proposition}

\begin{proof}
The first statement is the definition of geodesics of the rescaled metric, whose length functional is $C$. For constant $c$ it reduces to the length functional of $g$, whose geodesics are great circles. For the interface, take the plane approximation with a straight interface and endpoints at fixed distances from it; the cost of a path crossing at a point $y$ is $c_1\sqrt{a_1^2+y^2}+c_2\sqrt{a_2^2+(d-y)^2}$, and setting its derivative in $y$ to zero gives $c_1\sin\alpha_1=c_2\sin\alpha_2$.
\end{proof}

The refraction relation shows how the great circle is bent by terrain. If the cost per kilometer on one side of a ridge zone is three times that on the other, a path approaching the interface from the cheap side at $\alpha_1$ enters the costly zone at an angle $\alpha_2$ satisfying $\sin\alpha_2=\sin\alpha_1/3$, so that the path is steered toward the normal and crosses the costly zone by the shortest passage. The penalty of departing from the great circle is then $C(\gamma^\ast)-c\,R\theta\ge0$, which is the quantity to be reported for each arc in the assessment.
